% !TeX root = ../thesisv2.tex

\chapter{理论基础、数据集与实验设置}
\label{cha:v2-foundation}

\section{理论基础}
\subsection{卷积神经网络相关基础}
卷积神经网络通过局部感受野、权值共享和层级特征抽取机制，对输入图像进行逐层表达。对于语义分割任务而言，卷积层的作用并不仅仅是局部纹理检测，更重要的是通过连续卷积和下采样逐渐建立从边缘纹理到高级语义的特征层次。道路提取任务中，浅层卷积特征通常对边缘、亮度变化和局部线结构更敏感，而深层卷积特征则更适合表达道路与背景之间的语义差异。

在标准分类网络中，下采样通常以压缩空间分辨率为代价换取更强语义表达；而在语义分割任务中，模型既需要高层语义，也需要精确空间定位，因此编码器--解码器结构成为主流方案。编码器通过逐层下采样扩大感受野，解码器则逐步恢复空间分辨率。对于遥感道路提取而言，编码器负责识别道路走向、周边地物关系和区域级上下文，解码器负责恢复道路宽度、交叉口轮廓和支路边缘。只有二者协同工作，才能在复杂背景下得到结构完整的道路掩膜。

U-Net 的核心优势来自跳跃连接。跳跃连接直接将编码器浅层特征传递到解码器对应层，使深层语义与浅层边缘信息能够在恢复分辨率阶段完成融合\cite{unet2015}。这一结构对于道路提取尤为关键，因为道路目标通常较窄，且边缘信息决定了道路连通性和形态完整性。如果只依赖深层特征，细小支路和交叉口边界容易在下采样过程中被削弱；如果只依赖浅层特征，则又容易受到背景噪声影响。因此，跳跃连接的利用方式直接影响道路提取模型的最终效果。

卷积神经网络在训练时通常配合归一化、非线性激活和不同上采样方式使用。Batch Normalization 能够在一定程度上缓解训练中的分布漂移，提高梯度传播稳定性，但其效果与 batch size 有关；双线性上采样具有参数少、实现简单的特点；转置卷积则提供可学习的上采样能力。本文在不同模型中保留了这些结构差异，并通过统一整图评估观察其在道路提取任务中的作用。

\subsection{与本文相关的关键机制}
\subsubsection{跳跃连接特征重标定}
在 U-Net 中，跳跃连接承担浅层细节传递任务。由于遥感图像背景复杂，浅层特征中常同时包含道路边缘、建筑边界、阴影轮廓和地表纹理。因此，对跳跃连接特征进行重标定是一个具有明确意义的方向。本文主线模型 DLGU-Net 即围绕该问题展开，通过在四级跳跃连接处引入轻量级双线性注意力模块，对编码器特征在通道和空间两个维度同时进行筛选与增强。

从理论上看，通道重标定解决的是“哪一类特征更重要”的问题，空间重标定解决的是“特征图中哪些位置更重要”的问题。两者结合后，可以在不改变 U-Net 主干基本拓扑的前提下，提高解码器获得的浅层特征质量。这一路线与 Attention U-Net 的门控思想、UNet++ 的特征桥接思想具有共通的研究目标，即增强编码器与解码器之间的信息传递质量，但具体实现方式不同。

\subsubsection{多尺度上下文建模}
道路提取任务存在显著的尺度变化。主干道路通常宽度较大、走向明显，而支路、窄路和支线道路则更依赖局部细节判断。如果模型在高层特征中仅保留单一感受野信息，就可能出现主干道路提取稳定而细路漏检，或相反强调局部细节而背景误检偏高的问题。多尺度上下文建模正是为解决这一矛盾提出的策略。

空洞卷积通过在卷积核内部引入膨胀率，有效扩大等效感受野，同时保持卷积参数量在可控范围内\cite{dilated2016}。当膨胀率分别设置为 1、2、4 和 8 时，同一层并行分支可以分别观察局部区域、中等尺度区域和更大尺度区域特征。对于遥感道路提取任务而言，这意味着模型能够同时感知窄路边缘、主干道路走向以及道路与周边地物的空间关系。本文的空洞卷积增强 baseline 即围绕这一思想展开，在 bottleneck 位置引入并行多尺度上下文模块。

\subsubsection{整图级评估思想}
道路提取的最终使用对象是完整遥感图像中的道路网络，而不是独立的 patch。训练阶段出于显存和吞吐考虑通常采用 patch 输入，但如果直接使用 patch 级指标作为模型优劣结论，容易忽略道路的整图连通性、跨 patch 边界一致性和整体结构恢复质量。因此，本文全部核心结论统一建立在整图验证与整图测试协议之上。该协议先对整图进行补边和切块，再完成逐块预测和空间位置拼接，最后在整图层面统计指标。这样得到的 IoU 和 F1 更接近道路提取任务的真实使用场景。

\subsection{损失函数与评价指标}
本文主对比模型统一采用 Dice Loss 与二元交叉熵损失的组合形式。二元交叉熵从逐像素分类角度约束模型输出，使模型能够稳定区分道路与非道路像素；Dice Loss 直接关注预测区域与标注区域之间的重叠程度，能够在前景占比较低的情况下提升区域级优化效果。两者结合后，既兼顾数值稳定性，也兼顾区域一致性。

设模型对第 $i$ 个像素输出道路概率为 $p_i$，真实标签为 $y_i \in \{0,1\}$，则二元交叉熵损失可写为
\[
\mathcal{L}_{\mathrm{BCE}}=-\frac{1}{N}\sum_{i=1}^{N}\left[y_i\log p_i + (1-y_i)\log(1-p_i)\right].
\]
Dice 系数定义为
\[
\mathrm{Dice}=\frac{2\sum_i p_i y_i+\epsilon}{\sum_i p_i+\sum_i y_i+\epsilon},
\]
对应 Dice Loss 为
\[
\mathcal{L}_{\mathrm{Dice}}=1-\mathrm{Dice}.
\]
因此，本文组合损失写为
\[
\mathcal{L}=\mathcal{L}_{\mathrm{BCE}}+\mathcal{L}_{\mathrm{Dice}}.
\]

在评价指标方面，本文采用 IoU、F1、Precision、Recall 和 Accuracy 作为整图级评价指标。其中 IoU 作为主要排序依据，用于反映预测掩膜与真实掩膜之间的重叠程度；F1 反映 Precision 与 Recall 的综合平衡；Precision 体现模型对误报的控制能力；Recall 体现模型对漏检的控制能力；Accuracy 则提供总体像素分类正确率参考。设真阳性、假阳性、假阴性和真阴性分别为 $TP$、$FP$、$FN$ 和 $TN$，则有
\[
\mathrm{IoU}=\frac{TP}{TP+FP+FN},\qquad
\mathrm{Precision}=\frac{TP}{TP+FP},
\]
\[
\mathrm{Recall}=\frac{TP}{TP+FN},\qquad
\mathrm{F1}=\frac{2TP}{2TP+FP+FN}.
\]

\begin{table}[htbp]
\centering
\caption{本文主要评价指标的含义与分析侧重点}
\label{tab:v2-metric-meaning}
\begin{tabular}{p{2.2cm}p{4.0cm}p{6.0cm}}
\toprule
指标 & 数值含义 & 在本文中的主要分析作用 \\
\midrule
IoU & 预测区域与标注区域的交并比 & 全文主要排序指标，反映整体分割质量 \\
F1 & Precision 与 Recall 的调和平均 & 反映模型在误报与漏检之间的平衡能力 \\
Precision & 预测为道路的像素中真实为道路的比例 & 用于分析背景抑制能力和误报控制能力 \\
Recall & 真实道路像素中被正确预测的比例 & 用于分析细小道路、断裂道路和支路的保留能力 \\
Accuracy & 全体像素分类正确率 & 用于辅助观察模型整体分类稳定性 \\
\bottomrule
\end{tabular}
\end{table}

\section{数据集介绍}
\subsection{Massachusetts Roads Dataset}
当前本地材料中，全部完成训练、验证和测试的可复核实验均基于 Massachusetts Roads Dataset 展开。该数据集提供高分辨率遥感图像及其对应的道路标注，是道路提取研究中使用较广泛的公开数据集之一\cite{mnih2013}。由于该数据集具有较高分辨率和较为明确的道路标注，适合作为本文统一实验工程的核心数据来源。

在本文工程中，数据集被统一整理为训练集、验证集和测试集三个部分，并进一步区分为图像目录和标签目录，便于训练脚本、评估脚本和整图推理脚本直接调用。经过目录重组后，训练、验证和测试样本的组织方式完全一致，避免了原始下载来源差异对工程复现造成影响。

\begin{table}[htbp]
\centering
\caption{Massachusetts Roads Dataset 在本文工程中的数据规模}
\label{tab:v2-dataset-scale}
\begin{tabular}{lccc}
\toprule
数据划分 & 原始图像数 & 训练输入形式 & 最终评价形式 \\
\midrule
训练集 & 1108 & 离线 patch & patch 训练 \\
验证集 & 14 & 离线 patch + 整图评估 & 整图验证 \\
测试集 & 49 & 整图切块推理 & 整图测试 \\
\bottomrule
\end{tabular}
\end{table}

从数据规模可以看出，训练集数量明显大于验证集和测试集，这使本文在模型训练阶段能够积累较多样本，而在最终结果分析阶段则更适合采用整图级评价，以减小 patch 统计口径对结果解释造成的偏差。

\subsection{数据集样本特征}
图 \ref{fig:v2-dataset-examples} 展示了 Massachusetts Roads Dataset 测试集中的典型样例及其道路标注。可以看到，数据集同时包含道路结构相对简单的郊野区域、道路分布密度中等的城乡过渡区域，以及支路密集、背景纹理复杂的城镇区域。不同场景之间在道路密度、道路宽度、阴影分布和地表纹理方面差异明显，这为模型检验提供了较好的场景多样性。

从样本视觉特征分析，道路提取主要面临四类挑战。第一，主干道路与窄支路共存，尺度差异显著；第二，道路周围建筑和树木投影会破坏边界连续性；第三，停车场、裸土地表和建筑边缘与道路具有局部纹理相似性；第四，道路交叉口和分支区域对解码器细节恢复能力要求较高。这些数据特征决定了本文模型设计需要同时关注浅层边界恢复、背景抑制和多尺度上下文表达。

\begin{figure}[htbp]
\centering
\includegraphics[width=\textwidth]{generated/dataset_examples.png}
\caption{Massachusetts Roads Dataset 测试集样例及其道路标注}
\label{fig:v2-dataset-examples}
\end{figure}

\subsection{数据组织与离线 patch 构建}
原始遥感图像尺寸约为 $1500 \times 1500$。为了与网络输入尺寸保持一致，本文首先对图像和标签使用反射方式补边，将尺寸统一到 $1536 \times 1536$，随后将每张图像切分为 $3 \times 3$ 共 9 个 $512 \times 512$ patch。训练阶段直接读取离线 patch 文件，以减少重复整图解码和在线切块开销；验证与测试阶段则保留整图层面的切块预测和结果恢复流程。

离线 patch 数据集的建立是本文工程重建中的重要环节。早期实现中，模型每次读取样本都需要重新解码整张遥感图像、完成补边并从中裁出当前 patch，这种做法虽然逻辑清晰，但会使数据准备成为训练瓶颈。重构后，训练集共生成 9972 个 patch，验证集共生成 126 个 patch，模型训练阶段可以直接读取 patch 文件，从而显著提高数据吞吐效率。

\begin{figure}[htbp]
\centering
\includegraphics[width=0.95\textwidth]{generated/preprocessing_pipeline.png}
\caption{遥感道路样本的补边与 patch 划分流程}
\label{fig:v2-preprocess}
\end{figure}

\section{实验设置}
\subsection{实验环境与参数}
本文实验平台采用 Python 与 PyTorch 框架构建，主要依赖 OpenCV 进行图像读取与预处理，依赖 YAML 配置文件管理实验参数，依赖 Matplotlib 生成训练曲线和对比图。模型训练与测试均在 NVIDIA GeForce RTX 5090 显卡环境下完成。为了保证模型间的公平比较，主对比模型统一采用相同的输入尺寸、优化器、损失函数和学习率调度策略。

\begin{table}[htbp]
\centering
\caption{主对比模型统一训练配置}
\label{tab:v2-train-config}
\begin{tabular}{lc}
\toprule
配置项 & 数值 \\
\midrule
输入尺寸 & $512 \times 512$ \\
补边目标尺寸 & $1536 \times 1536$ \\
批大小 & 4 \\
初始学习率 & $1 \times 10^{-4}$ \\
权重衰减 & $1 \times 10^{-5}$ \\
优化器 & Adam \\
损失函数 & Dice + BCE \\
学习率调度器 & StepLR(10, 0.5) \\
训练轮数上限 & 20 \\
混合精度 & 开启 \\
梯度裁剪 & 1.0 \\
训练数据加载线程数 & 4 \\
验证数据加载线程数 & 0 \\
训练日志间隔 & 100 step \\
\bottomrule
\end{tabular}
\end{table}

上述配置具有两个作用。其一，它为全部模型提供统一训练条件，使模型性能差异主要由结构本身而非外部条件变化造成。其二，它在保证训练稳定性的同时控制了实验成本，使多组完整模型能够在同一硬件环境下顺利完成训练和测试。

\subsection{数据预处理}
本文的数据预处理流程包括图像读取、标签读取、尺寸统一、patch 划分、图像归一化和标签二值化等环节。图像统一按 RGB 顺序读取，并采用 ImageNet 均值和方差完成归一化，以提高模型训练稳定性。标签则统一处理为道路与非道路的二值图。训练阶段使用离线 patch；验证和测试阶段按照整图切块、拼接和裁剪恢复的流程进行。

这一预处理方式具备三个优点。第一，输入尺寸统一，有利于配置驱动训练与多模型共用数据管线。第二，离线 patch 将重复解码和补边操作移出训练环节，提高了 GPU 计算利用率。第三，整图恢复流程保证了最终评价对象与道路提取实际使用对象一致，从而使结果解释更具工程意义。

\subsection{评价指标}
虽然 IoU 和 F1 是本文最核心的两个排序指标，但在道路提取任务中，Precision 和 Recall 的配合同样重要。一个模型若 Precision 较高，通常说明其能够更有效地抑制建筑阴影、停车场纹理等非道路区域带来的误报；一个模型若 Recall 较高，则说明其对细小道路、断裂道路或支路的保留能力更强。本文在结果分析中始终同时给出 Precision 和 Recall，以避免仅凭单一指标解读模型行为。

Accuracy 在道路提取任务中提供总体像素分类正确率，但由于非道路像素占比通常较大，该指标并不适合作为排序主依据。因此，本文统一以 IoU 为第一排序指标，以 F1 为辅助排序指标，再通过 Precision 和 Recall 分析模型在背景抑制与支路保留之间的倾向。

\subsection{工程优化}
本文在原始实验实现基础上完成了数据管线和训练流程优化。核心改动包括：构建离线 patch 数据集、采用多进程 DataLoader、固定训练日志间隔、减少训练阶段频繁同步开销以及统一整图推理与整图评估逻辑。这些优化并不改变模型结构本身，但显著提升了实验效率和结果组织质量。

\begin{table}[htbp]
\centering
\caption{工程优化前后基线模型单 epoch 训练耗时}
\label{tab:v2-engineering-speed}
\begin{tabular}{lcc}
\toprule
训练阶段 & 单 epoch 时间/s & 说明 \\
\midrule
早期在线切块实现 & 409.3 & 在线整图解码、补边和切块 \\
离线 patch 优化后 & 119.1 & 离线 patch + 优化后的数据加载 \\
日志与同步进一步优化后 & 113.5 & 固定步长日志 + 减少同步开销 \\
\bottomrule
\end{tabular}
\end{table}

由表 \ref{tab:v2-engineering-speed} 可见，离线 patch 优化后，单 epoch 训练时间由 409.3 s 降至 119.1 s，训练耗时下降约 70.9\%；进一步结合日志与同步优化后，单 epoch 时间降至 113.5 s，相比早期实现总体下降约 72.3\%，对应约 3.61 倍的训练吞吐提升。该结果表明，数据管线优化是本文得以完成多模型完整实验的重要前提。

\begin{figure}[htbp]
\centering
\includegraphics[width=0.72\textwidth]{generated/engineering_speedup.png}
\caption{工程优化前后基线模型单 epoch 训练耗时对比}
\label{fig:v2-engineering-speedup}
\end{figure}

\subsection{整图评估协议}
本文的核心评估方式为整图评估协议。训练阶段使用 patch 数据进行梯度更新，验证与测试阶段则对完整遥感图像进行补边、切块预测、空间拼接和结果恢复，最终在整图层面统计 IoU、F1、Precision、Recall 和 Accuracy。这样做的主要原因在于，道路提取的最终目标是获得完整道路网络，而非若干独立 patch 的局部结论。

整图评估协议具有三方面意义。第一，它避免了 patch 边界效应对指标造成的直接干扰；第二，它使验证结果与测试结果采用同一统计口径，有利于后续模型选择和结论解释；第三，它更符合遥感道路提取的真实应用流程，能够直接体现道路网络的完整性、分支保持能力和背景抑制效果。因此，本文的全部模型排序、性能分析和结论均以整图评估为基础。

\section{本章小结}
本章围绕后续模型分析所依赖的共用条件，系统介绍了卷积神经网络与 U-Net 系列方法的相关理论基础、跳跃连接重标定和多尺度上下文建模等关键机制、Massachusetts Roads Dataset 的组织方式和样本特征，以及统一训练配置、数据预处理、工程优化和整图评估协议。通过离线 patch 数据集和统一整图验证流程，本文建立了适合多模型公平比较的实验基础，为后续主线模型和 baseline 改进分支的分析提供了统一前提。
